\scanpage{28}
\deleted{\unsure{(ii)|(iii)}}\inserted{(i)} $\Rightarrow(\mathrm{iii})$ \inserted{pour un $n$ donné,}
on note que (i) est stable par changement
plat de la base, et (iii) est \uline{invariant} par
changement \inserted{\uline{fid}} plat de la base. Cela nous ramène au
cas où les extensions résiduelles de $A'/A$ sont ttes
triviales. \deleted{\unsure{Mais alors on}} De plus, on peut évidemment
supposer \inserted{$A=A_n=B^{(n)}$, et} $A$ local [en le remplaçant au besoin
par un localisé]. Alors on voit que si
$E_{A'_n{}^{\deleted{\unsure{m}}}}$ \deleted{plat} sur $A'_n{}^{\deleted{\unsure{m}}}$ est muni d'une donnée de
descente, \deleted{\unsure{alors}} rel à $A'_n{}^{\deleted{\unsure{m}}}/A_n{}^{\deleted{\unsure{m}}}$
[i.e. à $A'_n/B^{(n)}$]
$E_{A'_n}$ est néc. un $A'_n$-module \uline{libre} [car il
suffit de voir que son rang \deleted{\unsure{aux}}\inserted{en les} différents pts
de $\operatorname{Spec} A'_n$ est le $\overline{\mathrm m}$, ce qui nous ramène au
cas où $A$ est le spectre d'un corps, où cela
résulte de l'effectivité de la donnée de descente.]
Je dis qu'on a effectivité pour $A'_n/B^{(n)}$, donc
un module descendu libre de $\overline{\mathrm m}$ rang sur $B^{(n)}$.
Cela résulte du lemme 1, cor. 3, compte tenu
du corollaire au lemme \unsure{?|2.}

On a ainsi prouvé a)

\textcircled{b} Pour prouver (i), comme la formation de
$H^1(A'/A)$ commute à l'extension plate de la base
$A$, on peut supposer $A$ \uline{complet}. Alors \deleted{\unsure{$B|E$}} on a
% archive sheet number: 27
